\documentclass[10pt,a4paper]{amsart}
\usepackage[margin=19mm]{geometry}
\usepackage{amsmath,amssymb,mathtools}
\usepackage[scr=rsfs,cal=euler]{mathalfa}
\usepackage{xfrac}
\usepackage[unicode,hidelinks,bookmarks=false]{hyperref}
\newcommand{\ie}{i.e.\ }
\newcommand{\cf}{cf.\ }
\newcommand{\resp}{respectively\ }
\newcommand{\Subsection}{\subsection}
\providecommand{\nobreakdash}{\nobreak\hbox{-}}
\setlength{\emergencystretch}{2em}
\multlinegap0pt
\usepackage{amssymb}
\usepackage{mathtools}
\usepackage{tikz}
\usepackage{algorithm}\usepackage{algpseudocode}
\newtheorem{theorem}{Theorem}
\title{Characterization of Blind Code Rate Recovery in Linear Block Codes}
\author{Atreya Vedantam, Radha Krishna Ganti}
\date{Source: arXiv:2603.02031v1 (CC BY 4.0)}
\begin{document}
\maketitle

\noindent\emph{Source and licence.} Characterization of Blind Code Rate Recovery in Linear Block Codes. Version arXiv:2603.02031v1 (CC BY 4.0), distributed by arXiv under the Creative Commons Attribution 4.0 International licence, http://creativecommons.org/licenses/by/4.0/, as recorded in the arXiv licence metadata for this exact version and shown on the abstract page https://arxiv.org/abs/2603.02031v1. Full licence notice: \texttt{licenses/arxiv-2603.02031-CC-BY-4.0.txt}.\par\smallskip

\noindent\textit{Editorial scope.} Counted unit: the article's theorem rankincrease (source0.tex lines 179--184). The article's complete original proof of it is delivered with it (source0.tex lines 185--188, one \texttt{proof} environment of 4 lines). No mathematics is written, repaired or re-derived: the statement and the proof are the article's own lines, in the article's own notation, and no retained line is substituted or re-flowed.  No further source-local context is needed: every cross-reference of every retained line resolves inside the two delivered spans.  Nothing was omitted from any retained span, so the package carries no omission record.  No retained line contains a citation, so the wrapper carries no bibliography and no bibliography entry.  No retained line contains a figure environment, an image-inclusion command or drawing code, so no image is delivered and the package declares no asset dependency.  Editorial formatting only, in addition: an AMS \texttt{amsart} preamble that restates the article's own declarations and macros used by the retained lines, namely the environments \texttt{theorem} on separate counters, together with the standard packages those lines need.  The retained environments print in this document as Theorem 1 in order of appearance, whereas the article prints the same items with the numbering of its own full document.  The complete native source of the article as distributed in the arXiv e-print for this exact version is bundled unchanged as \texttt{sources/195591/source0.tex}.
\begin{theorem}
    Suppose $p, m_1, m_2, \ldots, m_d \in \{0, 1\}^{M_s}$ are columns of the word matrix such that $p = \bigoplus_{i=1}^d m_i$. Let $p_e'$ be the probability of a bit error in the code word matrix. Suppose $M_s >> d$ and $p_e'$ was small such that $M_sp_e'$ is large. Then the probability that the rank increases by $1$ in the presence of at least one error (denoted by $p_1$) is lower bounded as follows.
    \begin{align}\label{rankincrease}
        p_1 \geq \Phi \left(\frac{(M_s-d)^2 - 2M_s^2(d+1)p_e'}{2M_s(d+1)\sqrt{M_sp_e'(1-p_e')}}\right).
    \end{align}
\end{theorem}

\begin{proof}
    See Appendix for a detailed proof. 
    Proof sketch: We (almost surely) lower bound the minimum distance of the code formed by the span of $\{ m_1, m_2, \ldots, m_d\}$. This distance increases with increasing $M_s$. Now given at least one bit error in the matrix, there must be enough errors to `cross' this minimum distance, if the linear dependency is to be retained. For a small bit error probability and high $M_s$, this becomes vanishingly unlikely, concluding the proof.
\end{proof}

\end{document}
